\begin{lemma}{Eigenschaften des Banach--Mazur--Abstands}
             {EigenschaftenDesBanachMazurAbstands}
Seien $(E, \norm\cdot_E), (F, \norm\cdot_F), (G, \norm\cdot_G)$
normierte $\K$--Vektorräume.\\
Dann gilt:
\begin{enumerate}[(1)]
	\item \label{EigenschaftenDesBanachMazurAbstands1}
	Sind $\Inv(E, F) \neq \emptyset$ und $\Inv(F, G) \neq \emptyset$,
	so auch $\Inv(E, G) \neq \emptyset$ und es gilt
	\[d_{BM}(E,G)\le d_{BM}(E,F)d_{BM}(F,G)\text.\]
	\item \label{EigenschaftenDesBanachMazurAbstands2}
	Sind $(E, \norm\cdot_E)$ und $(F, \norm\cdot_F)$ isometrisch isomorph, so ist
	$d_{BM}(E,F) = 1$.
	\item \label{EigenschaftenDesBanachMazurAbstands3}
	Sind $\norm\cdot$ eine weitere Norm auf $E$ und $c, d \in \R_{>0}$
	mit $c \norm\cdot \le \norm\cdot_E \le d \norm\cdot$, so ist
	\[d_{BM}\bigg((E, \norm\cdot_E), (E, \norm\cdot) \bigg) \le \frac d c
	\text.\]
	\item\label{EigenschaftenDesBanachMazurAbstands4}
	$d_{BM}(E,F)= d_{BM}(F,E)$.
\end{enumerate}
\end{lemma}
\begin{proof}
Es gilt für alle $T \in \Inv(E, F)$:
\[ 1 = \opnorm{\id_E}EE = \opnorm{T\inv \circ T} E E
\le \opnorm T EF \opnorm{T\inv}FE\]
und daher auch stets
\begin{align}d_{BM}\bigg((E, \norm\cdot_E), (F, \norm\cdot_F)\bigg) \ge 1\text.
\tag{$\ast$}\label{edbmaeq:1}
\end{align}
(1) Seien $R \in \Inv(E, F)$ und $S \in \Inv(F, G)$.\\
Dann ist $S \circ R$ linear, invertierbar und stetig. Wegen $(S \circ R)\inv
= R\inv \circ S\inv$ ist aber auch $(S \circ R)\inv$ stetig und linear und
damit $S \circ R \in \Inv(E, G)$.\\
Daher gilt auch
\begin{align*}d_{BM}\left( (E, \norm\cdot_E), (G, \norm\cdot_G)\right)
&\le \opnorm{S\circ R}EG \opnorm{(S \circ R)\inv}G E\\
&\le \opnorm SFG \opnorm{S\inv}GF \opnorm REF \opnorm{R\inv}FE\text.
\end{align*}
Wegen $\opnorm REF \opnorm{R\inv}FE > 0$ ist
$\frac{d_{BM}\left( (E, \norm\cdot_E), (G, \norm\cdot_G)\right)}
{\opnorm REF \opnorm{R\inv}FE}$ eine untere Schranke von
$\menge{\opnorm A F G \opnorm{A\inv}GF}{A \in \Inv(F,G)}$ und damit
\[\frac{d_{BM}\left( (E, \norm\cdot_E), (G, \norm\cdot_G)\right)}
{\opnorm REF \opnorm{R\inv}FE} \le 
d_{BM}\left( (F, \norm\cdot_F), (G, \norm\cdot_G)\right)\text.\]
Analog ist nun (wegen \eqref{edbmaeq:1} wohldef.) $\frac{d_{BM}\left((E, \norm\cdot_E), 
(G, \norm\cdot_G)\right)}{d_{BM}\left( (F, \norm\cdot_F), 
(G, \norm\cdot_G)\right)}$ eine untere Schranke von
\[\menge{\opnorm B E F \opnorm{B\inv}FE}{B \in \Inv(E,F)}\] und damit
\[\frac{d_{BM}\left( (E, \norm\cdot_E), 
(G, \norm\cdot_G)\right)}{d_{BM}\left( (F, \norm\cdot_F), 
(G, \norm\cdot_G)\right)} \le d_{BM}\left( (E, \norm\cdot_E), 
(F, \norm\cdot_F)\right)\text.\]
(2) Seien $(E, \norm\cdot_E)$ und $(F, \norm\cdot_F)$ isometrisch isomorph.\\
Dann existiert $\varphi : E \rightarrow F$ linear, stetig und bijektiv mit
$\norm \cdot _F \circ \varphi = \norm\cdot _E$.\\
Da $\varphi$ insbesondere eine bijektive Isometrie ist, ist auch $\varphi\inv$ stetig
und es gilt $\norm\cdot_F = \norm\cdot_E \circ \varphi \inv$.\\
Somit ist $\varphi \in \Inv(E, F)$ und es ist $\opnorm \varphi E F = 1$, sowie
$\opnorm {\varphi\inv} F E = 1$.
Also
\[d_{BM}\bigg((E, \norm\cdot_E), (F, \norm\cdot_F)\bigg)
\le \opnorm \varphi E F \opnorm {\varphi\inv} F E = 1\text.\]
Wegen \eqref{edbmaeq:1} gilt aber auch die umgekehrte Ungleichung, also
Gleichheit.\\
(3) Seien $\norm\cdot$ eine Norm auf $E$ und $c, d \in\!\ ]0,\infty[$
mit $c \norm\cdot \le \norm\cdot_E \le d \norm\cdot$.\\
Sei $I : (E, \norm\cdot) \rightarrow (E, \norm\cdot_E), x \mapsto x$.\\
Offenbar ist $I$ linear und invertierbar.\\
Wegen $\norm\cdot_E \le d \norm\cdot$ ist $I$ stetig und es gilt
$\opnorm I {(E, \norm\cdot)}{(E, \norm\cdot_E)} \le d$.\\
Analog ist wegen $\norm\cdot \le \frac 1 c \norm \cdot_E$ auch $I\inv$ stetig
und es gilt $\opnorm {I\inv}{(E, \norm\cdot_E)}{(E, \norm\cdot)}\le\frac 1c$.\\
Also ist
\[d_{BM}\bigg((E, \norm\cdot_E), (E, \norm\cdot) \bigg) \le
\opnorm I {(E, \norm\cdot)}{(E, \norm\cdot_E)}
\opnorm {I\inv}{(E, \norm\cdot_E)}{(E, \norm\cdot)} \le \frac d c\text.\]
(4) Offenbar gilt für $T \in L(E, F)$
\[T \in \Inv(E, F) \iff T\inv \in \Inv(F, E)\]
und $(T\inv)\inv = T$, sodass gilt:
\[\mengea{\opnorm S E F \opnorm{S\inv}FE}{S \in \Inv(E, F)}
= \mengea{\opnorm R F E \opnorm{R\inv}FE}{R \in \Inv(F, E)}\text.\]
Daher stimmen die Infima dieser Mengen überein.
\end{proof}